\section{Causal Validation (Companion Work)}\label{causal-validation-companion-work}

\label{app:causal}

The geometry reported in the main text is read off the
representation: it shows where foundation information is decodable.
Whether the model \emph{uses} that information during generation is a
causal question, taken up in full by companion work
\citep{reblitzrichardson2026causal}. Here we report the initial
causal checks on the OLMo-2\textasciitilde7B foundation directions used throughout
this paper, which indicate that the directions are functionally
implicated rather than probing artifacts.

\textbf{Direction ablation.} Projecting a foundation's direction out of the
residual stream at a target layer degrades that foundation's
continuations more than the other foundations'. Specificity, the gap
between the on-target and off-target mean effect over 48 prompts and
the six foundations, is negative at every layer tested and deepens
with depth, from \(-0.12\) at layer 4 to \(-0.32\) at layer 14. Removing a
foundation's direction selectively harms that foundation, and the
effect is strongest where the directions are most stable.

\textbf{Steering injection.} Adding \(\alpha\) times a foundation's direction
to the residual stream produces a dose-response. Mean specificity
rises monotonically with the injection strength, from \(+0.08\) at
\(\alpha = 1\) to \(+0.16\), \(+0.61\), \(+1.85\), and \(+3.59\) at
\(\alpha = 2, 5, 10, 20\). The same direction that decodes a foundation
also steers generation toward it, in proportion to the dose.

These results are descriptive rather than a full causal localization;
the companion paper \citep{reblitzrichardson2026causal} reports the
complete treatment, including behavioral grounding and sparse
autoencoder feature overlap. They establish that the foundation
directions in \Cref{results} correspond to features the model acts on during
generation, not only to features that are linearly readable.
